\documentclass[a4paper]{article}
\usepackage[a4paper, margin=1cm]{geometry}
\usepackage[utf8]{inputenc}
\usepackage[T1]{fontenc}
\usepackage{tikz}
\usetikzlibrary{shapes, arrows, fit, calc, positioning}
\usepackage{tikz-uml}
\tikzumlset{actor below=0.4cm, fill usecase=white, fill class=white, text=black, draw=black}
\pagestyle{empty}

\begin{document}

\begin{center}
{\Large Diagramma dei casi d'uso -- HomeStock}

\vspace{0.6cm}
\resizebox{\textwidth}{!}{%
\begin{tikzpicture}

\begin{umlsystem}[x=0, y=0]{HomeStock}
  % casi d'uso principali
  \umlusecase[x=0, y=0,     name=reg,   width=4cm]{Registrazione e accesso}
  \umlusecase[x=0, y=-1.8,  name=man,   width=4cm]{Aggiunta manuale prodotti}
  \umlusecase[x=0, y=-3.6,  name=exp,   width=4cm]{Date di scadenza}
  \umlusecase[x=0, y=-5.4,  name=inv,   width=4cm]{Visualizzazione inventario}
  \umlusecase[x=0, y=-7.2,  name=qta,   width=4cm]{Gestione quantità}
  \umlusecase[x=0, y=-9.0,  name=spesa, width=4cm]{Lista della spesa}
  \umlusecase[x=0, y=-10.8, name=ric,   width=4cm]{Suggerimento ricette}
  \umlusecase[x=0, y=-12.6, name=cond,  width=4cm]{Condivisione familiare}
  \umlusecase[x=0, y=-14.4, name=loc,   width=4cm]{Organizzazione per luogo}
  \umlusecase[x=0, y=-16.2, name=prem,  width=4cm]{Premium}
  \umlusecase[x=0, y=-18.0, name=stat,  width=4cm]{Statistiche e spreco}
  % casi d'uso raggiunti tramite include ed extend
  \umlusecase[x=10, y=-1.8,  name=scan,  width=4cm]{Scansione codice a barre}
  \umlusecase[x=10, y=-3.6,  name=notif, width=4cm]{Avvisi e notifiche}
  \umlusecase[x=10, y=-5.4,  name=cerca, width=4cm]{Ricerca e ordinamento}
  \umlusecase[x=10, y=-9.0,  name=trasf, width=4cm]{Trasferimento in dispensa}
  \umlusecase[x=10, y=-10.8, name=ingr,  width=5.2cm]{Ingredienti disponibili e mancanti}
  \umlusecase[x=10, y=-12.6, name=prio,  width=4cm]{Priorità ai prodotti in scadenza}
\end{umlsystem}

% attori
\umlactor[x=-8, y=0.2]{Visitatore}
\umlactor[x=-8, y=-5.6]{Membro famiglia}
\umlactor[x=-8, y=-14.6]{Utente}
\umlactor[x=-8, y=-19.6]{Utente Premium}

% generalizzazione tra attori:
% l'Utente (proprietario della casa) fa tutto quello che fa un Membro famiglia, piu' altro;
% l'Utente Premium fa tutto quello che fa un Utente, piu' altro
\umlinherit{Utente}{Membro famiglia}
\umlinherit{Utente Premium}{Utente}

% associazioni (disegnate sotto le ellissi)
\begin{pgfonlayer}{background}
\umlassoc[anchor2=180]{Visitatore}{reg}

\umlassoc[anchor2=180]{Membro famiglia}{reg}
\umlassoc[anchor2=180]{Membro famiglia}{man}
\umlassoc[anchor2=180]{Membro famiglia}{exp}
\umlassoc[anchor2=180]{Membro famiglia}{inv}
\umlassoc[anchor2=180]{Membro famiglia}{qta}
\umlassoc[anchor2=180]{Membro famiglia}{spesa}
\umlassoc[anchor2=180]{Membro famiglia}{ric}
\umlassoc[anchor2=180]{Membro famiglia}{cond}

\umlassoc[anchor2=180]{Utente}{loc}
\umlassoc[anchor2=180]{Utente}{prem}

\umlassoc[anchor2=180]{Utente Premium}{stat}
\end{pgfonlayer}

% extend: dall'estensione verso il caso base
\umlextend{scan}{man}
\umlextend{notif}{exp}
\umlextend{cerca}{inv}
\umlextend{trasf}{spesa}

% include: dal caso base verso il caso incluso
\umlinclude{ric}{ingr}
\umlinclude{ric}{prio}

\end{tikzpicture}%
}
\end{center}

\end{document}
